\begin{table}[!htbp]
\centering
\caption{基准返航安全余量下的最大安全载荷}
\label{tab:q1-payload}
\small
\setlength{\tabcolsep}{4pt}
\renewcommand{\arraystretch}{1.10}
\begin{tabular}{@{}lrrrrr@{}}
\toprule
服务区 & 距离（km） & 巡航海拔（m） & A（kg） & B（kg） & C（kg） \\
\midrule
S001 & 3.05 & 273.06 & 25.00 & 30.00 & 80.00 \\
S002 & 7.41 & 308.95 & 25.00 & 30.00 & 68.86 \\
S003 & 7.26 & 561.80 & 25.00 & 30.00 & 68.21 \\
S004 & 7.71 & 398.24 & 25.00 & 30.00 & 63.70 \\
S005 & 5.28 & 326.41 & 25.00 & 30.00 & 80.00 \\
S006 & 3.19 & 366.10 & 25.00 & 30.00 & 80.00 \\
S007 & 4.53 & 511.97 & 25.00 & 30.00 & 80.00 \\
S008 & 8.08 & 372.91 & 25.00 & 28.80 & 58.90 \\
S009 & 5.71 & 293.23 & 25.00 & 30.00 & 80.00 \\
S010 & 4.78 & 435.64 & 25.00 & 30.00 & 80.00 \\
S011 & 2.81 & 295.36 & 25.00 & 30.00 & 80.00 \\
S012 & 7.39 & 413.30 & 25.00 & 30.00 & 68.12 \\
S013 & 4.85 & 399.18 & 25.00 & 30.00 & 80.00 \\
S014 & 5.42 & 591.95 & 25.00 & 30.00 & 80.00 \\
S015 & 6.01 & 571.80 & 25.00 & 30.00 & 80.00 \\
\bottomrule
\end{tabular}
\par\smallskip{\footnotesize 最大载荷仅表示质量与能量允许的连续上界；实际组批还须满足体积约束。}
\end{table}
